\section{Discussion}
In this study, we integrated municipal street-tree inventories, PlanetScope-derived phenology, and 1-km near-surface urban air-temperature data to examine individual-tree phenology across 76 cities in the contiguous United States. Our results support three main conclusions. First, phenological variation within cities was often comparable to variation among cities for the same species, indicating that urban tree phenology cannot be fully represented by city-level averages. Second, phenological timing and growing season length were widely associated with fine-scale intra-city temperature variation, although the direction and magnitude of these associations differed among species, cities, and phenological phases. Third, phenological sensitivity to temperature varied systematically with regional climatic context: warmer local urban environments were generally associated with earlier spring onset and delayed fall senescence in cooler or wetter cities, whereas these associations weakened or reversed in hotter or drier cities. Together, these findings show that urban tree phenology reflects interactions among local thermal heterogeneity, species identity, and regional climate, with implications for urban forest monitoring, species selection, and climate adaptation.

\subsection{Fine-scale urban temperature and intra-city variation in phenology}

Our results show that urban phenology cannot be treated as a uniform city-level average (Questions 1 and 2). For many species, individual-tree phenological variation within cities was comparable to, or larger than, variation among cities (Fig.~\ref{fig1:variance} and Fig.~\ref{fig2:coefficients}). This finding is consistent with previous city-specific case studies in the U.S., showing that temperature regulates urban phenology through local micrometeorological conditions, despite differences in data sources \citep{Crawford2024ERL, Katz2019STOTEN, Parece2018UrbanSci, Wu2025NatCommun, Zipper2016ERL}. Our study advances this literature by revealing the broader generality of phenological heterogeneity and fine-scale temperature sensitivity across species and cities: these patterns are evident across a broad set of urban tree inventory species, and remain apparent across cities spanning diverse climatic regions (Fig.~\ref{fig2:coefficients} and Fig.~\ref{fig3:examples}). The individual-level framework captures species-specific and within-city spatial details that are often unresolved by coarser-resolution sensors \citep{badeck_responses_2004, zeng_review_2020, Zhao2022RSE}. 

These results also highlight the importance of fine-scale urban thermal heterogeneity. Previous studies have often examined urban phenology through urban--rural contrasts, macroclimatic gradients, or proxies for urbanization such as impervious surface fraction, land surface temperature, or population density \citep{Dallimer2016EcolEvol, jia_urbanization_2021, Li2017GCB, Meng2020PNAS, Wang2025NatCities, Wang2019NatEcoEvol}. By contrast, our analysis links individual-tree phenology to near-surface air temperature within cities, allowing us to evaluate phenological sensitivity to temperature across many species--city combinations. Because air temperature is a direct environmental cue for plant phenology, this approach provides a closer representation of the thermal conditions relevant to trees than many indirect urbanization metrics \citep{chakraborty_lower_2020, du_simultaneous_2021}. 

The observed intra-city heterogeneity has practical implications for urban forest management. In many species--city combinations, earlier SOS and later EOS under warmer local conditions were accompanied by longer GSL (Fig.~\ref{fig2:coefficients}b, d), indicating that fine-scale urban thermal variation can alter not only the timing of individual phenological events but also the duration of seasonal canopy activity. By altering the timing and magnitude of canopy-mediated functions, particularly seasonal carbon uptake and release, variation in individual-tree phenology may scale up to neighborhood-level differences in ecosystem services and disservices, including cooling benefits and pollen exposure \citep{Crawford2024ERL, Katz2019STOTEN, Song2025SRS}. Management strategies based only on citywide averages may overlook these fine-scale differences. Incorporating high-resolution monitoring of both urban thermal environments and tree phenology could improve urban forest planning under continued warming \citep{Hall2016LandscapeEcol}.

\subsection{Regional climate modulates phenological sensitivity to urban temperature}

Phenological sensitivity to temperature varied not only among species and cities, but also across regional climatic gradients. Spring phenology showed weaker and more variable associations with local temperature than fall phenology, whereas fall senescence was more consistently delayed in warmer local urban environments (Fig.~\ref{fig2:coefficients}, Fig.~\ref{fig4:map}a, b). This difference between SOS and EOS is consistent with broader phenological literature showing that spring and fall phenology are regulated by different combinations of environmental cues and physiological constraints. Spring leaf-out depends on interactions among winter chilling, spring forcing, photoperiod, and species-specific requirements \citep{Ettinger2020NCC, Fu2015Nature, Hanninen2019TPS, Li2022JAG, Meng2020PNAS}, whereas fall senescence can be strongly influenced by late-season temperature, water availability, carbon balance, and stress events \citep{Bigler2021AFM, Chaves1991JXB, Fu2018GCB, Luo2024FunctEcol}.

The climate-dependent patterns observed here suggest that local warming within cities does not have uniform phenological impacts. In cooler cities, warmer local urban environments were generally associated with earlier spring onset. In hotter cities, however, this association weakened or became positive, indicating delayed SOS in warmer local environments (Fig.~\ref{fig2:coefficients}a, Fig.~\ref{fig4:map}a, c). This pattern is consistent with the possibility that additional warming in already warm environments may reduce chilling accumulation or push trees closer to thermal stress thresholds, thereby limiting or reversing spring advancement \citep{Liang2024UFUG, Yin2024EarthsFuture}. However, because we did not directly quantify chilling fulfillment, forcing accumulation, or tree physiological status, this mechanism needs a direct test in future studies.

For fall phenology, warmer local urban environments were generally associated with delayed senescence in wetter cities but earlier senescence in drier cities (Fig.~\ref{fig2:coefficients}b, Fig.~\ref{fig4:map}b, d). This pattern is consistent with moisture limitation modifying the effect of local warming on fall phenology. In wetter environments, warming may extend photosynthetic activity and delay senescence when water limitation is weak. In drier environments, however, warmer local conditions may increase evaporative demand, intensify drought stress, or coincide with more impervious and sparsely vegetated microsites, leading to earlier senescence \citep{Franceschi2023UFUG, Liang2024UFUG, Stratopoulos2019IJB, Yin2024EarthsFuture}. Because soil moisture, vapor pressure deficit, and irrigation were not directly measured, these explanations should be viewed as plausible mechanisms for future testing.

These findings complement previous studies showing that urban phenological responses depend on regional climate \citep{Li2019NatEcoEvol, Meng2020PNAS, Yang2023GCB}. However, our results extend this perspective from urban--rural or city-level contrasts to individual trees within cities. The same degree of local urban warming may be associated with different phenological outcomes depending on whether a city is cool or warm, wet or dry. As a result, models that assume uniform phenological responses to urban warming may misestimate growing-season length, canopy function, and species suitability in urban forests \citep{Fu2023GCB, Li2019NatEcoEvol}.

\subsection{Species-specific phenological sensitivity and its potential drivers}

Species identity was an important axis of phenological heterogeneity for urban trees (Figs.~\ref{fig2:coefficients} and \ref{fig4:map}a,b). Even within the same city, co-occurring species often differed in both phenological timing and phenological sensitivity to temperature. Similarly, the same species sometimes showed different temperature sensitivities across cities. These patterns indicate that urban phenology reflects not only local environmental conditions, but also the taxonomic composition of planted urban forests.

Species-specific responses are expected because tree species differ in phenological cues, chilling and forcing requirements, drought tolerance, hydraulic strategies, photoperiod sensitivity, and stress tolerance \citep{cole_shifting_2017, Wu2025NatCommun}. In urban settings, these intrinsic differences may be amplified by planting practices. Many urban tree species are planted outside their native climatic ranges or in environments that differ from those experienced by source populations \citep{Chamberlain2023CCE, ghelardini_variation_2006, Jensen2022EcolAppl, linkosalo_twilight_2006}. Provenance and nursery origin may therefore influence phenological responses, because phenology is often heritable and locally adapted \citep{Chamberlain2023CCE, Iler2021AREES, Wu2025NatCommun}. For example, southern-origin populations often show delayed budburst in cooler or higher-latitude urban environments because of differences in photoperiod sensitivity, chilling requirements, and temperature thresholds \citep{ghelardini_variation_2006, Jensen2022EcolAppl, linkosalo_twilight_2006}, while northern provenances may require fewer growing degree days and thus respond more readily to urban warming \citep{Chamberlain2023CCE}. Future work incorporating detailed climatic legacy, planting history, and physiological processes will be needed to disentangle these effects.

These species-level differences have important implications for urban forest planning. As cities warm, species that maintain appropriate phenological timing under local thermal and moisture conditions may be better able to sustain canopy function, cooling, carbon uptake, and biotic interactions \citep{Dallimer2016EcolEvol, Ettinger2022NCC,Harrison2015FunctEcol}. Conversely, species showing delayed spring onset or earlier fall senescence in warmer or drier local environments may be more vulnerable to functional stress, especially if these phenological shifts reflect unmet chilling requirements, heat stress, or water limitation in urban environments \citep{Du2022JAG,Zani2020Science}. Phenological sensitivity to temperature may therefore provide an early indicator of species performance under urban climate stress, complementing occurrence-based or demographic assessments of urban tree suitability \citep{Burley2019STOTEN,EsperonRodriguez2022NCC,Wu2025NatCommun,Zhu2025UrbanEcosyst}.

\subsection{Limitations and future directions}

Several limitations should be considered when interpreting our results. First, although PlanetScope provides frequent observations at relatively high spatial resolution, satellite-derived phenology for individual urban trees remains uncertain. Extracted EVI time series may include signals from neighboring crowns, understory vegetation, turfgrass, impervious surfaces, or other nearby land-cover features. This issue is likely more pronounced for small trees, such as those with canopy diameters smaller than 3 m, whose crowns may not be fully resolved at PlanetScope resolution \citep{Song2025SRS, Wu2021ISPRS}. Our validation against NPN observations, together with previous evaluation analyses, showed that PlanetScope-derived phenology captured meaningful seasonal signals, but agreement differed between SOS and EOS \citep{Song2025SRS, Zhao2022RSE}. Future work using standardized field observations, drone imagery, or high-resolution aerial data could better quantify uncertainty in individual-tree phenology retrieval.

Second, municipal tree inventories differ in timing, completeness, taxonomic resolution, and location accuracy \citep{McCoy2022eLife}. Some inventories may predate the PlanetScope observation period, and tree removals, new plantings, or changes in health condition may not be fully captured. Species misidentification or coordinate error could also contribute to apparent intra-city phenological variation. In addition, our inferences apply primarily to inventoried deciduous street trees, not all urban trees. Park trees, private-yard trees, and remnant urban forest patches may show different phenological patterns because they differ in canopy structure, soil conditions, management, and disturbance history.

Third, although U-HAT provides high-resolution near-surface air temperature across cities, its 1-km resolution does not capture all tree-scale thermal variation. Street trees can experience substantial microclimatic differences at meter- to tens-of-meters scales due to shading, building geometry, pavement, traffic, irrigation, and local canopy cover. Therefore, our analysis characterizes fine-scale intra-urban temperature variation at the neighborhood scale rather than the full microclimatic environment experienced by each tree. Integrating sub-kilometer air-temperature observations, mobile sensing, or sensor networks would improve the ability to estimate tree-scale thermal exposure.

Fourth, our seasonal temperature metrics summarize broad preseason conditions and may not fully capture the phenological cues that regulate SOS and EOS. Spring phenology is shaped by the interaction between winter chilling and spring forcing, while fall phenology may respond to late-season heat, cold spells, drought stress, and short-term weather variability \citep{Chuine2010Science, Fu2015Nature, wang_overestimation_2020}. Future analyses using chilling accumulation, growing degree days, vapor pressure deficit, soil moisture, and event-based temperature metrics could improve mechanistic attribution.

Finally, other urban environmental factors may contribute to intra-city phenological heterogeneity. Artificial light at night, in particular, can delay fall senescence and modify phenological sensitivity to temperature and precipitation \citep{Chen2026NatCommun, Friulla2025UrbanSci, Meng2022PNASNexus, Wang2025NatCities}. Land-use composition, irrigation, soil compaction, pollution, and management intensity may also interact with temperature to influence phenological timing. Integrating these variables with individual-tree phenology and urban climate data would provide a more complete understanding of the environmental drivers of urban phenological variation.

\subsection{Conclusion}

In this study, we showed widespread intra-city phenological variation and fine-scale temperature-phenology associations across 76 cities in the contiguous United States. Furthermore, we found the magnitude and direction of these associations varied substantially among species, and were modulated by regional climatic conditions. Recognizing that urban tree phenology emerges from interactions among local thermal heterogeneity, species identity, and regional climate can support both fine-scale monitoring of phenological responses and more informed species selection under continued urban warming. Our findings also suggest the need to move beyond single-driver explanations and more fully investigate the complex, context-dependent processes that shape urban tree phenology across spatial scales.